%% ============================================================
%%  PREAMBLE — Production-Grade Agentic AI: LangGraph + Python
%%  Extracted from main.tex.  Input ONCE from Tex/main.tex.
%% ============================================================

\documentclass[
  11pt,
  twoside,
  openright
]{book}

\usepackage{polyglossia}
\setdefaultlanguage{english}

\usepackage{fontspec}
\setmainfont{TeX Gyre Termes}[
  BoldFont   = TeX Gyre Termes Bold,
  ItalicFont = TeX Gyre Termes Italic
]
\setsansfont{TeX Gyre Heros}[
  BoldFont   = TeX Gyre Heros Bold,
  ItalicFont = TeX Gyre Heros Italic
]
\setmonofont{TeX Gyre Cursor}[
  Scale    = 0.88,
  BoldFont = TeX Gyre Cursor Bold
]

\usepackage{geometry}
\geometry{
  papersize  = {7in, 10in},
  top        = 1in,
  bottom     = 1.1in,
  inner      = 1in,
  outer      = 0.875in,
  headheight = 23.07pt,
  twoside
}

%% Fix fancyhdr \headheight warning (must be after geometry)
\setlength{\headheight}{23.07004pt}
\addtolength{\topmargin}{-0.00003pt}

\usepackage{xcolor}
\definecolor{BookBlue}       {HTML}{1A4A8A}
\definecolor{BookDark}       {HTML}{1C1C1E}
\definecolor{BookMid}        {HTML}{4A4A4A}
\definecolor{LearnBg}        {HTML}{EEF7EE}
\definecolor{LearnBorder}    {HTML}{27AE60}
\definecolor{TipBg}          {HTML}{EAF4FB}
\definecolor{TipBorder}      {HTML}{2980B9}
\definecolor{WarningBg}      {HTML}{FEF9E7}
\definecolor{WarningBorder}  {HTML}{F39C12}
\definecolor{NoteBg}         {HTML}{F4F4F4}
\definecolor{NoteBorder}     {HTML}{888888}
\definecolor{TakeawayBg}     {HTML}{F0F0FF}
\definecolor{TakeawayBorder} {HTML}{5B4FBE}
\definecolor{CodeBg}         {HTML}{F7F7F7}
\definecolor{CodeComment}    {HTML}{6A737D}
\definecolor{CodeKeyword}    {HTML}{D73A49}
\definecolor{CodeString}     {HTML}{032F62}
\definecolor{CodeConstant}   {HTML}{005CC5}
\definecolor{CodeAnnotation} {HTML}{6A737D}

\usepackage{microtype}
\usepackage{setspace}
\setstretch{1.18}

\usepackage{csquotes}
\usepackage{parskip}
\setlength{\parskip}{6pt}
\setlength{\parindent}{0pt}

\widowpenalty=10000
\clubpenalty=10000
\displaywidowpenalty=10000
\predisplaypenalty=10000
\raggedbottom
\emergencystretch=3em
\hbadness=3000
\vbadness=3000

\usepackage{needspace}
\usepackage{etoolbox}
\preto\section{\needspace{5\baselineskip}}
\preto\subsection{\needspace{4\baselineskip}}
\preto\subsubsection{\needspace{3\baselineskip}}

\usepackage{titlesec}
\renewcommand{\bottomtitlespace}{0.15\textheight}

\titleformat{\part}[display]
  {\centering\sffamily\Huge\bfseries\color{BookBlue}}
  {\sffamily\large\mdseries\partname\ \thepart}{20pt}{}

\titleformat{\chapter}[display]
  {\sffamily\huge\bfseries\color{BookBlue}}
  {\sffamily\normalsize\mdseries\chaptertitlename\ \thechapter}{12pt}{}
\titlespacing*{\chapter}{0pt}{0pt}{24pt}

\titleformat{\section}
  {\sffamily\Large\bfseries\color{BookBlue}}
  {\thesection}{1em}{}
\titlespacing*{\section}{0pt}{18pt}{6pt}

\titleformat{\subsection}
  {\sffamily\large\bfseries\color{BookDark}}
  {\thesubsection}{1em}{}
\titlespacing*{\subsection}{0pt}{14pt}{4pt}

\titleformat{\subsubsection}
  {\sffamily\normalsize\bfseries\color{BookDark}}
  {\thesubsubsection}{1em}{}
\titlespacing*{\subsubsection}{0pt}{10pt}{2pt}

\usepackage{fancyhdr}
\pagestyle{fancy}
\fancyhf{}
\fancyhead[LE]{\small\sffamily\color{BookMid}\leftmark}
\fancyhead[RO]{\small\sffamily\color{BookMid}\rightmark}
\fancyfoot[LE,RO]{\small\sffamily\color{BookMid}\thepage}
\renewcommand{\headrulewidth}{0.3pt}
\renewcommand{\footrulewidth}{0pt}

\fancypagestyle{plain}{
  \fancyhf{}
  \fancyfoot[LE,RO]{\small\sffamily\color{BookMid}\thepage}
  \renewcommand{\headrulewidth}{0pt}
}

\makeatletter
\renewcommand{\cleardoublepage}{%
  \clearpage
  \if@twoside
    \ifodd\c@page\else
      \thispagestyle{empty}\hbox{}\newpage
      \if@twocolumn\hbox{}\newpage\fi
    \fi
  \fi
}
\makeatother

\usepackage{rotating}
\usepackage{booktabs}
\usepackage{tabularx}
\usepackage{array}
\usepackage{longtable}
\renewcommand{\arraystretch}{1.35}

\newcolumntype{P}[1]{>{\RaggedRight\arraybackslash}p{#1}}

\usepackage{graphicx}
\usepackage{float}
\usepackage{caption}
\captionsetup{
  font      = {small,sf},
  labelfont = {bf,color=BookBlue},
  skip      = 6pt
}

\usepackage{listings}
\lstset{
  backgroundcolor   = \color{CodeBg},
  basicstyle        = \ttfamily\fontsize{9}{11}\selectfont,
  keywordstyle      = \color{CodeKeyword}\bfseries,
  stringstyle       = \color{CodeString}\itshape,
  commentstyle      = \color{CodeComment}\itshape,
  numbers           = left,
  numberstyle       = \tiny\color{BookMid},
  numbersep         = 6pt,
  breaklines        = true,
  frame             = single,
  framerule         = 0.4pt,
  rulecolor         = \color{NoteBorder},
  showstringspaces  = false,
  tabsize           = 2,
  captionpos        = b,
  breakatwhitespace = false,
  columns           = fullflexible,
  linewidth         = \linewidth,
  postbreak         = \mbox{\textcolor{gray}{$\hookrightarrow$}\space},
  xleftmargin       = 1.5em,
  framexleftmargin  = 1.5em,
  aboveskip         = 6pt,
  belowskip         = 4pt,
}

\lstdefinelanguage{Python}{
  morekeywords={
    False,None,True,and,as,assert,async,await,break,class,continue,def,
    del,elif,else,except,finally,for,from,global,if,import,in,is,lambda,
    nonlocal,not,or,pass,raise,return,try,while,with,yield
  },
  morekeywords=[2]{self,cls},
  sensitive=true,
  morecomment=[l]{\#},
  morestring=[b]",
  morestring=[b]',
  morestring=[s]{"""}{"""}, 
  morestring=[s]{'''},
  keywordstyle=\color{CodeKeyword}\bfseries,
  keywordstyle=[2]\color{CodeConstant},
  stringstyle=\color{CodeString},
  commentstyle=\color{CodeComment}\itshape,
  breaklines=true,
  breakatwhitespace=false,
  columns=fullflexible,
  keepspaces=true,
  showstringspaces=false,
  postbreak=\mbox{\textcolor{gray}{$\hookrightarrow$}\space},
}

\lstdefinelanguage{yaml}{
  morecomment=[l]{\#},
  morestring=[b]",
  morestring=[b]',
  literate=
    *{true}{{\color{CodeKeyword}true}}{4}
    {false}{{\color{CodeKeyword}false}}{5}
    {null}{{\color{CodeKeyword}null}}{4},
  sensitive=true,
  keywordstyle=\color{CodeKeyword}\bfseries,
  stringstyle=\color{CodeString},
  commentstyle=\color{CodeComment}\itshape,
  breaklines=true,
  columns=fullflexible,
  showstringspaces=false
}

\lstdefinelanguage{json}{
  morestring=[b]",
  morecomment=[l]//,
  morecomment=[s]{/*}{*/},
  sensitive=true,
  stringstyle=\color{CodeString},
  commentstyle=\color{CodeComment}\itshape,
}

\DeclareRobustCommand{\code}[1]{\texttt{#1}}
\providecommand{\term}[1]{\textit{#1}}
\providecommand{\chapterquote}[2]{%
  \vspace{1em}
  \begin{quote}\itshape#1\end{quote}
  \begin{flushright}--- #2\end{flushright}
  \vspace{1em}
}

\newcommand{\githubref}[1]{%
  \vspace{2pt}%
  {\small\sffamily\color{BookMid}Full source: \texttt{#1}}%
  \vspace{4pt}%
}

\usepackage{tcolorbox}
\tcbuselibrary{skins,breakable}

\newtcolorbox{learningbox}{
  colback=LearnBg, colframe=LearnBorder,
  before skip=12pt, fonttitle=\sffamily\bfseries\small,
  title={Learning Objective}, breakable, lines before break=3,
  left=8pt, right=8pt, top=6pt, bottom=6pt,
  boxrule=0.6pt, width=\linewidth
}
\newtcolorbox{tipbox}{
  colback=TipBg, colframe=TipBorder,
  before skip=12pt, fonttitle=\sffamily\bfseries\small,
  title={Tip}, breakable, lines before break=3,
  left=8pt, right=8pt, top=6pt, bottom=6pt,
  boxrule=0.6pt, width=\linewidth
}
\newtcolorbox{warningbox}{
  colback=WarningBg, colframe=WarningBorder,
  before skip=12pt, fonttitle=\sffamily\bfseries\small,
  title={Warning}, breakable, lines before break=3,
  left=8pt, right=8pt, top=6pt, bottom=6pt,
  boxrule=0.6pt, width=\linewidth
}
\newtcolorbox{notebox}{
  colback=NoteBg, colframe=NoteBorder,
  before skip=12pt, fonttitle=\sffamily\bfseries\small,
  title={Note}, breakable, lines before break=3,
  left=8pt, right=8pt, top=6pt, bottom=6pt,
  boxrule=0.6pt, width=\linewidth
}
\newtcolorbox{takeawaybox}[1][Key Takeaway]{
  colback=TakeawayBg, colframe=TakeawayBorder,
  before skip=12pt, fonttitle=\sffamily\bfseries\small,
  title={#1}, breakable, lines before break=3,
  left=8pt, right=8pt, top=6pt, bottom=6pt,
  boxrule=0.6pt, width=\linewidth
}

%% ── Source-Layout tree environment ─────────────────────────────────────
%% Usage:
%%   \begin{sourcelayout}
%%     \dirtree{%
%%       .1 src/.
%%       .2 agent/.
%%       .3 nodes/.
%%       .4 perceive\_node.py \DTcomment{§1.2}.
%%     }
%%   \end{sourcelayout}
%%
%% Requires: \usepackage{dirtree}  (loaded below)
%% Colours reuse NoteBg / NoteBorder (neutral grey — reference material,
%% not instructional content).
%% ────────────────────────────────────────────────────────────────────────
\usepackage{dirtree}

\newtcolorbox{sourcelayout}{%
  colback      = NoteBg,
  colframe     = NoteBorder,
  colbacktitle = NoteBorder,
  coltitle     = white,
  fonttitle    = \sffamily\bfseries\small,
  title        = {Source Layout},
  before skip  = 12pt,
  after skip   = 10pt,
  breakable,
  lines before break = 3,
  left   = 6pt,
  right  = 6pt,
  top    = 6pt,
  bottom = 6pt,
  boxrule = 0.6pt,
  width   = \linewidth,
}

%% dirtree global style: monospaced font, consistent with code listings
\renewcommand*\DTstylecomment{\rmfamily\color{CodeComment}\itshape}
\renewcommand*\DTstyle{\ttfamily\color{BookDark}}

\usepackage{amsmath}
\usepackage{amssymb}

\PassOptionsToPackage{hyphens}{url}
\usepackage{hyperref}
\usepackage{imakeidx}
\makeindex[columns=2, title=Index, intoc=true]
\urlstyle{tt}
\makeatletter
\g@addto@macro\UrlBreaks{\do\.\do\_\do\-}
\makeatother

\hypersetup{
  colorlinks  = true,
  linkcolor   = BookBlue,
  urlcolor    = BookBlue,
  citecolor   = BookBlue,
  pdftitle    = {Production-Grade Agentic AI: LangGraph + Python Frameworks},
  pdfauthor   = {Haythem Aber},
  pdfsubject  = {LangGraph, Python, Agentic AI, Production Systems},
  pdfkeywords = {LangGraph, Python, AI agents, production, observability, security}
}

\usepackage{tikz}
\usetikzlibrary{arrows.meta, positioning, shapes.geometric, shapes.misc, calc}
\usepackage{pdflscape}
\usepackage{ragged2e}

\newcommand{\component}[1]{\textsf{\textbf{#1}}}
\newcommand{\seechapter}[1]{\textit{(See Chapter~\ref{#1})}}
\newcommand{\seepart}[1]{\textit{(See Part~\ref{#1})}}

\newcommand{\partclose}[1]{%
  \cleardoublepage
  \vspace*{\fill}
  \begin{center}\sffamily\Large\color{BookBlue}#1\end{center}
  \vspace*{\fill}
  \cleardoublepage
}
